\section{Death and the Intermediate States}

Death is not the end of mind in this model. It is a structured passage, and for a trained practitioner it is the single greatest opportunity for liberation. Tibetan Buddhism maps dying as a \textbf{staged dissolution}. The gross body breaks down, then the subtle mind unfolds through fixed signs, ending in the clear light of death.

The completion-stage yogas and the bardo teachings are rehearsals for this passage. To die well is to recognize, at each stage, what is happening.

\subsection{The dissolution of the elements}

As the gross body dies, the \textbf{five elements dissolve in order}, each producing an outer sign, a vision seen by the dying person. Each dissolution is also the failing of one sense and one aggregate.

\begin{longtable}{l L{4.5cm} L{4.5cm}}
\caption{The outer dissolution and its four signs in order.}\\
\toprule
\textbf{Step} & \textbf{Element dissolves} & \textbf{Outer sign} \\
\midrule
1 & earth into water & mirage, a shimmering haze \\
2 & water into fire & smoke, smoky wisps \\
3 & fire into wind & fireflies, red sparks \\
4 & wind into consciousness & a guttering lamp, then a steady flame \\
\bottomrule
\end{longtable}

The signs come in fixed order. \textbf{Mirage}, \textbf{smoke}, \textbf{fireflies}, \textbf{lamp}. As earth dissolves the body feels heavy and sinking. As water dissolves the fluids dry. As fire dissolves the body heat leaves. As wind dissolves the outer breath stops. The person now looks dead. But the inner dissolution is just beginning.

\subsection{The white, the red, and the black}

After outer breath stops, the subtle minds dissolve through three colored visions. The two drops in the central channel release.

\begin{longtable}{l L{4.0cm} L{5.0cm}}
\caption{The inner dissolution, the white, the red, and the black.}\\
\toprule
\textbf{Step} & \textbf{Sign} & \textbf{What happens} \\
\midrule
1 & white appearance & the white drop descends from the crown, a vision like white moonlight, the thought-states of anger cease \\
2 & red increase & the red drop rises from the navel, a vision like red sunlight, the thought-states of attachment cease \\
3 & black near-attainment & white and red meet at the heart, a vision of black darkness, then unconsciousness, the thought-states of ignorance cease \\
\bottomrule
\end{longtable}

\subsection{The clear light of death}

After the blackness, the \textbf{clear light of death} (prabhasvara) dawns. This is the ground luminosity, the mother clear light, the very subtle mind laid bare. \textbf{This is the actual moment of death and the supreme chance for liberation.}

A trained yogi rests in the clear light and can attain liberation or the illusory body here. An untrained person passes through it unaware in an instant. The white, red, black, and clear light are the \textbf{four empties}.

\subsection{The three bardos}

If there is no recognition of the clear light, consciousness moves through the \textbf{intermediate states} (\textbf{bardos}).

\begin{longtable}{L{3.2cm} L{8.0cm}}
\caption{The three bardos of death.}\\
\toprule
\textbf{Bardo} & \textbf{What it is} \\
\midrule
bardo of dying & the death process itself, the dissolutions ending in the clear light \\
bardo of dharmata & the luminous visionary state, where the peaceful and wrathful deities appear \\
bardo of becoming & the drift toward rebirth, the mental body wandering, lasting up to forty-nine days \\
\bottomrule
\end{longtable}

The \textbf{Bardo Thodol}, known in English as the Tibetan Book of the Dead, is read aloud to guide the dead person through these three.

\subsection{The hundred deities as the mind's own nature}

In the \textbf{bardo of dharmata}, the visionary state, the \textbf{hundred peaceful and wrathful deities} appear. They are \textbf{not external gods}. They are the display of one's own awareness.

The count splits into \textbf{forty-two peaceful deities}, residing in the heart center, who appear first, and \textbf{fifty-eight wrathful deities}, residing in the crown center, who appear later to those who did not recognize the peaceful ones. Forty-two plus fifty-eight is one hundred.

The peaceful deities are the pure nature of the five poisons, the five wisdoms in serene form. The wrathful deities are the same energy shown as force, to cut through fixation. The whole set is generated from the \textbf{five buddha families}. Everything in the vision is a permutation of the five.

\subsection{Recognition as liberation}

The single point of the system is this. The deities, the lights, and the visions are nothing but the mind's own nature appearing to itself. \textbf{To recognize them as one's own awareness is liberation.} To fail to recognize them, to flee the wrathful forms and grasp at the dim lights of the six realms, is to be drawn back into rebirth. The death teaching and the bardo teachings exist to train this one recognition in advance, so it can be done when it counts.
